\documentclass[../DoAn.tex]{subfiles}
\begin{document}

Chương~\ref{chapter:Related_works} đã xác lập các yêu cầu cho một hệ thống gợi ý nhạc Việt dựa trên cảm xúc, không dùng dữ liệu người dùng và có cơ sở khoa học. Chương này trình bày các nền tảng lý thuyết và công nghệ được lựa chọn để đáp ứng các yêu cầu đó. Với mỗi thành phần, báo cáo nêu rõ nó giải quyết yêu cầu nào, các lựa chọn thay thế có thể dùng, và lý do của lựa chọn cuối cùng.

\section{Mô hình cảm xúc Valence--Arousal}
\label{section:3.1}
Để cả bài hát và màu sắc có thể được so khớp trên cùng một thước đo, đồ án cần một biểu diễn cảm xúc thống nhất. Mô hình vòng tròn cảm xúc của Russell \cite{russell1980circumplex} mô tả cảm xúc bằng hai trục trực giao: \textit{Valence} (mức độ dễ chịu, từ tiêu cực đến tích cực) và \textit{Arousal} (mức độ kích thích/năng lượng, từ trầm lắng đến sôi động); trục Arousal tương hợp với mô hình tâm trạng \textit{năng lượng--căng thẳng} của Thayer \cite{thayer1989biopsychology}, vốn nhấn mạnh chiều năng lượng của cảm xúc. Mỗi trạng thái cảm xúc trở thành một điểm trên mặt phẳng hai chiều này, và các nhãn cảm xúc rời rạc quen thuộc tương ứng với các vùng (góc phần tư) của mặt phẳng. Hình~\ref{fig:circumplex} minh họa bốn góc phần tư cùng tám nhãn cảm xúc mà hệ thống sử dụng.

\begin{figure}[tbp]
\centering
\begin{tikzpicture}[scale=1.0, font=\small,
   emo/.style={circle, fill=blue!55, inner sep=1.6pt}]
  % axes
  \draw[-{Stealth[length=2.5mm]}] (-3.4,0) -- (3.6,0) node[right] {Valence};
  \draw[-{Stealth[length=2.5mm]}] (0,-3.4) -- (0,3.6) node[above] {Arousal};
  \draw[dashed, gray] (0,0) circle (2.5);
  % quadrant labels
  \node[align=center, gray] at (1.9,3.1) {\textbf{Q1}};
  \node[align=center, gray] at (-1.9,3.1) {\textbf{Q2}};
  \node[align=center, gray] at (-1.9,-3.1) {\textbf{Q3}};
  \node[align=center, gray] at (1.9,-3.1) {\textbf{Q4}};
  % 8 emotions on the circle
  \node[emo] (e1) at (60:2.5) {}; \node[above right, font=\scriptsize] at (e1) {Phấn khích};
  \node[emo] (e2) at (25:2.5) {}; \node[right, font=\scriptsize] at (e2) {Vui vẻ};
  \node[emo] (e3) at (120:2.5) {}; \node[above left, font=\scriptsize] at (e3) {Căng thẳng};
  \node[emo] (e4) at (155:2.5) {}; \node[left, font=\scriptsize] at (e4) {Giận dữ};
  \node[emo] (e5) at (205:2.5) {}; \node[left, font=\scriptsize] at (e5) {U sầu};
  \node[emo] (e6) at (240:2.5) {}; \node[below left, font=\scriptsize] at (e6) {Buồn};
  \node[emo] (e7) at (300:2.5) {}; \node[below right, font=\scriptsize] at (e7) {Thư thái};
  \node[emo] (e8) at (335:2.5) {}; \node[right, font=\scriptsize] at (e8) {Bình yên};
\end{tikzpicture}
\caption{Mô hình vòng tròn cảm xúc của Russell: bốn góc phần tư (Q1 vui--sôi nổi, Q2 căng--giận, Q3 buồn--u sầu, Q4 thư thái--bình yên) và tám nhãn cảm xúc hệ thống sử dụng}
\label{fig:circumplex}
\end{figure}

Đây là mô hình được dùng phổ biến nhất trong nhận dạng cảm xúc âm nhạc và tâm lý học màu sắc, đồng thời là khung lý thuyết cho phép màu và nhạc gặp nhau, do liên kết màu--nhạc được trung gian hóa bởi cảm xúc \cite{palmer2013music}. So với các mô hình cảm xúc rời rạc (vui, buồn, giận...), mô hình hai trục liên tục cho phép tính khoảng cách và xếp hạng, nên phù hợp với bài toán truy hồi: ``gần nhau về cảm xúc'' trở thành ``gần nhau trên mặt phẳng V-A''. Mô hình hai trục có một hạn chế đã biết là gộp một số cảm xúc cùng vùng (ví dụ giận dữ và sợ hãi đều ở vùng Valence thấp--Arousal cao), điều mà các mô hình ba chiều (thêm trục Dominance) \cite{russellmehrabian1977pad} hay các thang đo cảm xúc chuyên biệt cho âm nhạc như GEMS \cite{zentner2008gems} xử lý tốt hơn; tuy nhiên với bài toán truy hồi và xếp hạng theo tâm trạng, sự phân biệt tinh tế này không thiết yếu, nên mô hình hai trục vẫn là lựa chọn phù hợp và phổ biến nhất. Vì vậy, toàn bộ hệ thống quy về không gian Valence--Arousal.

Có hai trường phái lớn trong biểu diễn cảm xúc. Trường phái \textit{rời rạc} (categorical) coi cảm xúc là một tập hữu hạn các trạng thái cơ bản (vui, buồn, giận, sợ...), tiêu biểu là sáu cảm xúc cơ bản của Ekman \cite{ekman1992basic}. Trường phái \textit{liên tục} (dimensional) như mô hình của Russell coi cảm xúc là điểm trong một không gian ít chiều liên tục. Với bài toán gợi ý, biểu diễn liên tục có hai lợi thế quyết định: nó cho phép định nghĩa \textit{khoảng cách} giữa hai trạng thái cảm xúc (để truy hồi ``gần nhau''), và cho phép \textit{nội suy} một đường chuyển cảm xúc mượt giữa hai điểm (phục vụ hành trình hai màu). Một biểu diễn rời rạc không có hai tính chất này: hai nhãn ``buồn'' và ``u sầu'' hoặc giống hệt hoặc khác hẳn, không có mức độ ở giữa. Trong hệ thống, các nhãn tâm trạng hiển thị cho người dùng (như ở Hình~\ref{fig:circumplex}) được suy ra \textit{từ} tọa độ V-A bằng cách xác định góc phần tư và vùng chứa điểm đó --- tức là nhãn rời rạc chỉ là một lớp trình bày trên nền biểu diễn liên tục, chứ không phải đơn vị tính toán.

Để làm rõ vì sao đồ án chọn mô hình hai trục của Russell thay vì các khung cảm xúc khác, Bảng~\ref{table:emotion-models} đối chiếu bốn lựa chọn tiêu biểu theo tiêu chí quan trọng nhất với bài toán gợi ý: liệu mô hình có cho phép định nghĩa \textit{khoảng cách} và \textit{nội suy} giữa các trạng thái hay không. Mô hình của Russell và mô hình PAD ba chiều \cite{russellmehrabian1977pad} đều là không gian liên tục có metric, nhưng trục Khống chế (Dominance) thứ ba của PAD rất khó suy ra một cách đáng tin từ âm thanh hay lời bài hát và thường cho dữ liệu thưa, nên lợi ích tách biệt được những cặp cảm xúc cùng vùng (giận dữ và sợ hãi) không bù lại chi phí ước lượng thêm một trục. Thang đo GEMS \cite{zentner2008gems} được thiết kế riêng cho cảm xúc \textit{do âm nhạc gợi lên} và tinh tế hơn về mặt mô tả, nhưng chín nhân tố của nó không tạo thành một không gian metric thuần nhất để tính khoảng cách màu--nhạc, đồng thời không có sẵn quy trình ánh xạ màu sang chín nhân tố này. Các mô hình rời rạc kiểu sáu cảm xúc cơ bản của Ekman \cite{ekman1992basic} trực giác và dễ gán nhãn nhưng thiếu hẳn hai tính chất cốt lõi mà truy hồi cần. Vì hai chức năng của đồ án đều quy về thao tác ``tìm điểm lân cận trên mặt phẳng cảm xúc'', một không gian liên tục hai chiều vừa đủ giàu để phân biệt tâm trạng vừa đủ gọn để ước lượng đáng tin là lựa chọn cân bằng nhất.

\begin{longtable}{|p{2.7cm}|p{2.0cm}|p{3.5cm}|p{3.6cm}|}
\caption{So sánh các mô hình biểu diễn cảm xúc và mức phù hợp với bài toán truy hồi}\label{table:emotion-models}\\
\hline
\textbf{Mô hình} & \textbf{Dạng} & \textbf{Ưu điểm} & \textbf{Hạn chế cho truy hồi} \\
\hline
\endfirsthead
\hline
\textbf{Mô hình} & \textbf{Dạng} & \textbf{Ưu điểm} & \textbf{Hạn chế cho truy hồi} \\
\hline
\endhead
Vòng tròn Russell (V-A) \cite{russell1980circumplex} & Liên tục, 2 chiều & Cho khoảng cách \& nội suy; phổ biến nhất trong MER và tâm lý học màu & Gộp một số cảm xúc cùng vùng (giận/sợ) \\
\hline
PAD \cite{russellmehrabian1977pad} & Liên tục, 3 chiều & Thêm trục Khống chế, tách được giận/sợ & Trục thứ ba khó suy từ âm thanh/lời; dữ liệu thưa \\
\hline
GEMS \cite{zentner2008gems} & 9 nhân tố nhạc & Thiết kế riêng cho cảm xúc do nhạc gợi & Không phải không gian metric; khó hợp nhất với màu \\
\hline
Cảm xúc cơ bản Ekman \cite{ekman1992basic} & Rời rạc & Trực giác, dễ gán nhãn & Không có khoảng cách/nội suy giữa các nhãn \\
\hline
\end{longtable}

Một điểm cần nhấn mạnh là tính \textit{trực giao} được giả định giữa hai trục: Valence và Arousal độc lập với nhau, nên một bài có thể vui--trầm lắng (Valence cao, Arousal thấp, như một bản acoustic nhẹ nhàng) hoặc buồn--dữ dội (Valence thấp, Arousal cao, như một bản rock bi tráng). Chính tính trực giao này cho phép bốn góc phần tư mô tả những kết hợp cảm xúc mà một thang đo một chiều (chỉ ``tích cực--tiêu cực'') không thể nắm bắt. Trong dữ liệu nhạc Việt thực tế của đồ án, hai trục có tương quan dương nhẹ (mục~\ref{subsection:eval-color}); đây là một thuộc tính của kho nhạc chứ không phải phản bác giả định lý thuyết, và nó có hệ quả thực tiễn là các vùng lệch trục (vui--trầm, buồn--dữ dội) thưa dữ liệu hơn.

Để chuyển từ tọa độ liên tục sang nhãn tâm trạng hiển thị cho người dùng, hệ thống xác định góc phần tư từ vị trí của điểm so với trung điểm $0.5$ trên mỗi trục, theo Phương trình~\eqref{eq:quadrant}, rồi gán nhãn cảm xúc theo vùng tương ứng (chẳng hạn Valence cao--Arousal cao thuộc góc ``vui vẻ/phấn khích''). Cần nhấn mạnh rằng đây thuần túy là một bước \textit{trình bày}: mọi phép truy hồi và xếp hạng vẫn thực hiện trên tọa độ liên tục, còn nhãn rời rạc chỉ sinh ra ở bước cuối để hiển thị.
\begin{equation}\label{eq:quadrant}
Q(v,a) = \begin{cases}
Q_1\ (\text{vui--sôi nổi}) & v \ge 0.5,\ a \ge 0.5 \\
Q_2\ (\text{căng--giận}) & v < 0.5,\ a \ge 0.5 \\
Q_3\ (\text{buồn--u sầu}) & v < 0.5,\ a < 0.5 \\
Q_4\ (\text{thư thái--bình yên}) & v \ge 0.5,\ a < 0.5
\end{cases}
\end{equation}

\section{Cơ sở liên kết màu sắc với cảm xúc}
\label{section:3.2}
Để ánh xạ màu về tọa độ cảm xúc, đồ án sử dụng ba cơ sở. Thứ nhất, để xử lý màu theo cảm nhận thị giác đồng đều, màu được chuyển sang không gian Oklab \cite{ottosson2020oklab} thay vì RGB thô. Không gian RGB không đồng đều về tri giác: cùng một khoảng cách số trong RGB có thể ứng với mức khác biệt cảm nhận rất khác nhau ở các vùng màu khác nhau. Oklab được thiết kế để khoảng cách Euclid xấp xỉ khác biệt cảm nhận, với trục $L$ (độ sáng) và hai trục màu $a$ (lục--đỏ), $b$ (lam--vàng). Từ tọa độ Oklab, đồ án suy ra hai đại lượng tri giác phục vụ ánh xạ cảm xúc: \textit{độ đỏ} (redness) tính từ góc màu $h$ và \textit{độ bão hòa} (saturation), theo Phương trình~\eqref{eq:redness}.

\begin{equation}\label{eq:redness}
  \text{redness} = \frac{1 + \cos h}{2}, \qquad h = \operatorname{atan2}(b, a)
\end{equation}

Thứ hai, trục Valence của màu được hồi quy về các chuẩn liên kết màu--cảm xúc của con người trong khảo sát ICEAS \cite{jonauskaite2020universal}, một khảo sát quy mô lớn trên nhiều quốc gia (Hình~\ref{fig:iceas-heatmap}). Thứ ba, trục Arousal của màu được suy ra theo quan hệ màu--âm nhạc của Whiteford \cite{whiteford2018color}: nhạc nhanh, sôi động gắn với màu sáng và bão hòa hơn. Đặc biệt, độ đỏ và độ bão hòa được đặt trên trục Arousal, phù hợp với phát hiện rằng độ bão hòa là yếu tố chính chi phối Arousal của màu \cite{valdez1994effects}. Hình~\ref{fig:oklab} minh họa hướng tăng của hai trục cảm xúc trong mặt phẳng màu. Lựa chọn dùng Whiteford thay vì chỉ dùng quan hệ màu đơn thuần xuất phát từ việc cần đúng cấu trúc màu--âm nhạc, vì hệ thống ánh xạ màu sang nhạc chứ không phải mô tả cảm nhận màu độc lập; công thức Arousal cụ thể được trình bày ở mục~\ref{subsection:4.2.2}.

Việc chọn không gian màu không phải là chi tiết kỹ thuật nhỏ. RGB là không gian thiết bị, không đồng đều về tri giác và trộn lẫn độ sáng với màu sắc. CIELAB cải thiện tính đồng đều nhưng vẫn có hiện tượng không đồng đều ở vùng xanh lam. Oklab \cite{ottosson2020oklab} là một cải tiến gần đây được thiết kế riêng để khoảng cách Euclid khớp tốt hơn với khác biệt cảm nhận, đồng thời tách bạch độ sáng ($L$) khỏi sắc độ ($a, b$) --- điều quan trọng vì hai trục cảm xúc của màu phụ thuộc vào các thành phần tri giác khác nhau (Valence gắn nhiều với độ sáng, Arousal gắn nhiều với độ bão hòa). Về phía cảm xúc, khảo sát ICEAS \cite{jonauskaite2020universal} thu thập liên kết màu--cảm xúc trên hàng nghìn người ở khoảng ba mươi quốc gia, cung cấp một chuẩn định lượng đáng tin cho trục Valence của màu. Riêng trục Arousal, ngoài Whiteford, phát hiện của Valdez và Mehrabian \cite{valdez1994effects} trên khoảng bảy mươi sáu màu Munsell cho thấy \textit{độ bão hòa} là yếu tố chi phối Arousal mạnh hơn cả độ sáng --- đây là cơ sở để công thức Arousal đặt trọng số đáng kể cho độ bão hòa.

Việc đặt màu và nhạc lên cùng một trục cảm xúc không phải là một quy ước tùy tiện mà có cơ sở trong hiện tượng \textit{tương ứng đa giác quan} (crossmodal correspondences): con người, một cách nhất quán và xuyên văn hóa, liên kết các thuộc tính của những giác quan khác nhau --- chẳng hạn âm vực cao với màu sáng, nhịp nhanh và âm lượng lớn với màu rực rỡ --- và phần lớn các liên kết này được \textit{trung gian hóa bởi cảm xúc} chứ không phải bởi tương đồng vật lý trực tiếp \cite{spence2011crossmodal, palmer2013music}. Nói cách khác, một màu và một bản nhạc ``hợp nhau'' khi chúng gợi lên cùng một trạng thái cảm xúc; do đó cảm xúc là cầu nối tự nhiên, và việc quy cả hai về tọa độ Valence--Arousal chỉ là hình thức hóa lại cầu nối sẵn có đó. Đây cũng là lý do đồ án không ghép màu với nhạc bằng quy tắc thủ công (ví dụ ``màu xanh dương luôn ứng với nhạc buồn'') mà ghép gián tiếp qua một không gian cảm xúc có cơ sở định lượng. Bảng~\ref{table:color-bases} tóm tắt ba cơ sở khoa học được dùng cho ánh xạ này và vai trò riêng của từng cơ sở.

\begin{longtable}{|p{3.0cm}|p{6.6cm}|p{2.4cm}|}
\caption{Ba cơ sở khoa học dùng cho ánh xạ màu sang tọa độ cảm xúc}\label{table:color-bases}\\
\hline
\textbf{Cơ sở} & \textbf{Vai trò trong ánh xạ} & \textbf{Nguồn} \\
\hline
\endfirsthead
\hline
\textbf{Cơ sở} & \textbf{Vai trò trong ánh xạ} & \textbf{Nguồn} \\
\hline
\endhead
Không gian Oklab & Xử lý màu đồng đều về tri giác; tách độ sáng ($L$) khỏi sắc độ ($a,b$) để hai trục cảm xúc dựa trên các thành phần khác nhau & \cite{ottosson2020oklab} \\
\hline
Chuẩn ICEAS & Hồi quy trục Valence của màu theo liên kết màu--cảm xúc do con người chấm (gần 30 quốc gia) & \cite{jonauskaite2020universal} \\
\hline
Quan hệ Whiteford & Suy trục Arousal theo quan hệ màu--\textit{âm nhạc} (màu sáng, bão hòa hơn ứng với nhạc nhanh hơn) & \cite{whiteford2018color} \\
\hline
Phát hiện Valdez--Mehrabian & Củng cố vai trò của độ bão hòa là yếu tố chi phối Arousal của màu & \cite{valdez1994effects} \\
\hline
\end{longtable}

\begin{figure}[tbp]
\centering
\includegraphics[width=0.72\textwidth]{Hinhve/iceas_heatmap.png}
\caption{Bản đồ nhiệt liên kết màu--cảm xúc từ khảo sát ICEAS trên khoảng ba mươi quốc gia: mỗi ô là xác suất Bayes liên kết một cảm xúc (hàng) với một màu (cột); ô càng cam/đỏ đậm thì liên kết càng mạnh (Nguồn: Jonauskaite và cộng sự, 2020 \cite{jonauskaite2020universal})}
\label{fig:iceas-heatmap}
\end{figure}

\begin{figure}[tbp]
\centering
\begin{tikzpicture}[scale=1.0, font=\small]
  \draw[-{Stealth[length=2.5mm]}] (-3.0,0) -- (3.2,0) node[right] {$a$ (lục $\rightarrow$ đỏ)};
  \draw[-{Stealth[length=2.5mm]}] (0,-3.0) -- (0,3.2) node[above] {$b$ (lam $\rightarrow$ vàng)};
  \draw[dashed, gray] (0,0) circle (2.4);
  \node[font=\scriptsize, gray] at (2.55,-0.3) {đỏ};
  \node[font=\scriptsize, gray] at (-2.5,-0.3) {lục};
  \node[font=\scriptsize, gray] at (0.4,2.5) {vàng};
  \node[font=\scriptsize, gray] at (0.4,-2.5) {lam};
  % arousal direction (toward red/saturated)
  \draw[-{Stealth[length=2.5mm]}, very thick, red!70] (0,0) -- (35:2.3)
       node[pos=1.18, align=center, font=\scriptsize] {Arousal $\uparrow$\\(đỏ, bão hòa, sáng)};
  % valence note
  \node[align=center, font=\scriptsize, blue!60!black] at (-1.5,1.7)
       {Valence $\uparrow$:\\sáng hơn ($L$),\\ít đỏ hơn};
\end{tikzpicture}
\caption{Mặt phẳng màu Oklab ($a$--$b$) và hướng tăng của hai trục cảm xúc: Arousal tăng theo độ đỏ và độ bão hòa (Whiteford), Valence tăng theo độ sáng và giảm theo độ đỏ (ICEAS)}
\label{fig:oklab}
\end{figure}

Việc tách bạch độ sáng khỏi sắc độ trong Oklab có một hệ quả thiết kế quan trọng: hai trục cảm xúc của màu dựa trên những thành phần tri giác \textit{khác nhau} và do đó tương đối độc lập. Trục Valence chủ yếu gắn với độ sáng $L$ (màu sáng hơn dễ chịu hơn) và mức độ đỏ, trong khi trục Arousal chủ yếu gắn với độ bão hòa và độ đỏ (màu rực rỡ kích thích hơn). Nếu dùng RGB thô --- nơi độ sáng bị trộn lẫn vào cả ba kênh --- thì một thay đổi nhằm điều chỉnh Valence sẽ vô tình kéo theo Arousal, làm hai trục dính vào nhau. Oklab gỡ được sự dính này, cho phép hồi quy mỗi trục cảm xúc trên đúng các thành phần tri giác chi phối nó. Một trường hợp cần xử lý riêng là các màu \textit{vô sắc} (trắng, xám, đen) có độ bão hòa gần không: với chúng, độ đỏ và sắc độ không còn ý nghĩa, nên Valence và Arousal được suy chủ yếu từ độ sáng (màu càng sáng càng dễ chịu, càng tối càng trầm lắng). Cách phân nhánh này bảo đảm ánh xạ màu vẫn hợp lý ở các đầu mút của dải màu.

\section{Biểu diễn âm nhạc bằng học tự giám sát}
\label{section:3.3}
Để đo độ ``nghe giống nhau'' và để suy ra Arousal từ âm thanh, đồ án cần một biểu diễn âm nhạc mạnh. Các đặc trưng âm thanh thủ công (như MFCC, chroma) nắm bắt cảm xúc yếu, nên đồ án dùng mô hình học tự giám sát (Self-Supervised Learning, SSL). Học tự giám sát huấn luyện mô hình trên lượng lớn nhạc \textit{không nhãn} bằng các bài toán tự sinh (dự đoán đoạn bị che, lượng tử hóa đặc trưng), nhờ đó thu được biểu diễn tổng quát mà không cần nhãn cảm xúc. Hai ứng viên hàng đầu là MERT \cite{li2023mert}, dùng cơ chế thầy--trò với mục tiêu âm học và âm nhạc, và MuQ \cite{zhu2025muq}, dùng lượng tử hóa vector dư trên phổ Mel; cả hai đều được đánh giá trên bộ chuẩn MARBLE \cite{yuan2023marble}.

Cơ chế của học tự giám sát có thể hình dung như sau: mô hình được cho ``nghe'' một đoạn nhạc đã bị che hoặc làm nhiễu một phần, rồi phải dự đoán phần bị thiếu dựa trên ngữ cảnh xung quanh. Để làm tốt việc này, mô hình buộc phải học các quy luật nội tại của âm nhạc --- hòa âm, nhịp điệu, âm sắc, cấu trúc --- mà không cần bất kỳ nhãn nào do con người gán. Sau khi tiền huấn luyện trên hàng nghìn giờ nhạc, mỗi lớp biến đổi (transformer) của mô hình mã hóa thông tin ở một mức trừu tượng khác nhau; đồ án gộp trung bình (mean-pooling) đầu ra theo trục thời gian thành một vector 1024 chiều duy nhất cho mỗi bài, đủ cô đọng để so khớp nhanh nhưng vẫn giàu thông tin. Chính vì biểu diễn này được học từ quy luật âm nhạc \textit{tổng quát} chứ không phải từ một tác vụ cảm xúc cụ thể, nó chuyển giao tốt sang nhạc Việt dù mô hình chưa từng được huấn luyện riêng cho nhạc Việt --- một tính chất then chốt với bài toán thiếu dữ liệu bản địa.

Hình~\ref{fig:ssl} mô tả cách một mô hình SSL âm nhạc được sử dụng trong đồ án. Tín hiệu âm thanh đi qua khối tiền xử lý, qua nhiều lớp biến đổi (transformer) đã được \textit{đóng băng}, rồi được gộp trung bình thành một vector biểu diễn 1024 chiều. Vector này phục vụ ba vai trò: (i) đo tương đồng giữa hai bài bằng cosine theo Phương trình~\eqref{eq:cosine}; (ii) bảo đảm nhất quán âm thanh cho gợi ý theo màu; và (iii) làm đặc trưng đầu vào cho đầu dò tuyến tính suy ra cảm xúc (mục~\ref{section:3.4}).

\begin{figure}[tbp]
\centering
\resizebox{\textwidth}{!}{%
\begin{tikzpicture}[font=\small,
   blk/.style={draw, rounded corners, align=center, minimum height=0.9cm, fill=black!3},
   arr/.style={-{Stealth[length=2mm]}, thick}]
  \node[blk, text width=1.7cm] (wav) at (0,0) {Tín hiệu\\âm thanh};
  \node[blk, text width=1.9cm] (pre) at (2.6,0) {Tiền xử lý\\(Mel / patch)};
  \node[blk, text width=2.3cm, fill=blue!6] (tr) at (5.6,0) {Transformer\\(đóng băng) $\times N$};
  \node[blk, text width=1.7cm] (pool) at (8.5,0) {Gộp\\trung bình};
  \node[blk, text width=1.9cm, fill=green!7] (emb) at (10.9,0) {Vector\\1024 chiều};
  \draw[arr] (wav)--(pre); \draw[arr] (pre)--(tr); \draw[arr] (tr)--(pool); \draw[arr] (pool)--(emb);
  % hai nhánh sử dụng, rẽ từ một điểm ngay dưới vector embedding (không chồng mũi tên)
  \node[blk, text width=3.3cm] (sim) at (8.7,-2.3) {Cosine $\rightarrow$ độ tương đồng};
  \node[blk, text width=3.3cm] (probe) at (12.7,-2.3) {Đầu dò tuyến tính $\rightarrow$ Arousal};
  \coordinate (fork) at (10.9,-1.15);
  \draw[arr] (emb) -- (fork) -- (sim.north);
  \draw[arr] (fork) -- (probe.north);
\end{tikzpicture}}
\caption{Sơ đồ khối sử dụng mô hình học tự giám sát âm nhạc: từ tín hiệu âm thanh tới vector biểu diễn 1024 chiều, dùng cho đo tương đồng và làm đặc trưng cho đầu dò cảm xúc}
\label{fig:ssl}
\end{figure}

\begin{equation}\label{eq:cosine}
  \cos(\mathbf{x},\mathbf{y}) = \frac{\mathbf{x}^\top \mathbf{y}}{\lVert \mathbf{x}\rVert \, \lVert \mathbf{y}\rVert}
\end{equation}

Hai mô hình khác nhau ở \textit{bài toán tự giám sát} (pretext task). MERT học bằng cách dự đoán các nhãn giả sinh từ một ``thầy âm học'' (gom cụm đặc trưng âm thanh) và một ``thầy âm nhạc'' (biểu diễn dựa trên cao độ), nhờ đó nắm được cả kết cấu âm sắc lẫn thông tin hòa âm. MuQ thay thế nhãn giả bằng mục tiêu lượng tử hóa vector dư trên phổ Mel (Mel Residual Vector Quantization), cho mục tiêu học ổn định hơn (Hình~\ref{fig:muq-arch}). Bộ chuẩn MARBLE \cite{yuan2023marble} đánh giá các mô hình này trên nhiều tác vụ hiểu nhạc (gắn thẻ, nhận dạng nhạc cụ, phát hiện cảm xúc...). Một điểm cốt lõi trong thiết kế của đồ án là dùng các mô hình này ở dạng \textit{đóng băng}: trọng số của mô hình không được cập nhật, chỉ một lớp tuyến tính phía trên được huấn luyện. Cách này phù hợp với bằng chứng rằng đầu dò tuyến tính trên đặc trưng đóng băng tổng quát hóa tốt hơn tinh chỉnh khi dữ liệu đích lệch miền \cite{kumar2022finetuning}, và đặc biệt cần thiết khi dữ liệu nhãn cảm xúc nhạc Việt khan hiếm.

Lý do dùng học tự giám sát thay vì đặc trưng âm thanh thủ công (MFCC, chroma, các bộ mô tả phổ) nằm ở chỗ các đặc trưng thủ công được thiết kế để mô tả tín hiệu chứ không để nắm bắt cảm xúc hay phong cách, nên thường yếu khi dùng làm đầu vào cho nhận dạng cảm xúc âm nhạc \cite{yang2012mer}. Học tự giám sát, ngược lại, học một biểu diễn tổng quát từ hàng nghìn giờ nhạc bằng các bài toán tự sinh (ví dụ che một phần tín hiệu rồi yêu cầu mô hình khôi phục), nhờ đó vector thu được mã hóa được cả âm sắc, hòa âm lẫn cấu trúc --- những yếu tố mà cảm xúc và độ ``nghe giống'' phụ thuộc vào --- mà không cần một nhãn cảm xúc nào trong quá trình tiền huấn luyện. Đây chính là tính chất quyết định với bài toán nhạc Việt thiếu nhãn.

Bảng~\ref{table:muq-mert} đối chiếu hai ứng viên theo các tiêu chí then chốt. Khác biệt cốt lõi nằm ở \textit{bài toán tự giám sát}: MERT phải dựa vào hai ``thầy'' ngoài (một thầy âm học gom cụm đặc trưng, một thầy âm nhạc dựa trên cao độ) để sinh nhãn giả, trong khi MuQ tự sinh mục tiêu bằng lượng tử hóa vector dư trên phổ Mel, cho quá trình học ổn định hơn và không phụ thuộc chất lượng của thầy ngoài. Trên bộ chuẩn MARBLE \cite{yuan2023marble}, MuQ ngang hoặc nhỉnh hơn MERT ở nhiều tác vụ hiểu nhạc; tuy nhiên đồ án không lấy thứ hạng trên bộ chuẩn làm căn cứ chấp nhận mà kiểm lại bằng độ đo cuối của chính hệ thống (mục~\ref{subsection:eval-bakeoff}).

\begin{longtable}{|p{2.9cm}|p{4.5cm}|p{4.5cm}|}
\caption{So sánh hai mô hình học tự giám sát âm nhạc MERT và MuQ}\label{table:muq-mert}\\
\hline
\textbf{Tiêu chí} & \textbf{MERT} \cite{li2023mert} & \textbf{MuQ} \cite{zhu2025muq} \\
\hline
\endfirsthead
\hline
\textbf{Tiêu chí} & \textbf{MERT} \cite{li2023mert} & \textbf{MuQ} \cite{zhu2025muq} \\
\hline
\endhead
Bài toán tự giám sát & Dự đoán nhãn giả từ thầy âm học + thầy âm nhạc (cao độ) & Lượng tử hóa vector dư trên phổ Mel (Mel-RVQ) \\
\hline
Tính chất mục tiêu học & Cần thầy ngoài; nhiều nhãn giả & Tự sinh mục tiêu, học ổn định hơn \\
\hline
Đánh giá trên MARBLE & Mạnh trên nhiều tác vụ hiểu nhạc & Ngang hoặc hơn MERT \cite{yuan2023marble} \\
\hline
Vai trò trong đồ án & Backbone tiền nhiệm (để so sánh) & Backbone được chọn cho cả ba vai trò \\
\hline
Phán quyết theo độ đo cuối & --- & Ngang/hơn; thắng chuyển giao GPT (mục~\ref{subsection:eval-bakeoff}) \\
\hline
\end{longtable}

Đồ án lựa chọn MuQ làm trục biểu diễn âm thanh sau một so sánh công bằng cho thấy MuQ \textit{ngang hoặc hơn} MERT --- hòa trên các độ đo trực tiếp và vượt có ý nghĩa thống kê ở phép chuyển giao sang một bộ chấm (judge) độc lập (chi tiết ở mục~\ref{subsection:eval-bakeoff}). Biểu diễn của MuQ được dùng cho cả ba vai trò nêu trên, bảo đảm một trục biểu diễn nhất quán cho toàn hệ thống.

\begin{figure}[tbp]
\centering
\includegraphics[width=0.6\textwidth]{Hinhve/muq_arch.png}
\caption{Kiến trúc mô hình MuQ: phổ Mel bị che một phần được Conformer dự đoán lại các token Mel-RVQ qua mục tiêu lượng tử hóa vector dư (Nguồn: Zhu và cộng sự, 2025 \cite{zhu2025muq})}
\label{fig:muq-arch}
\end{figure}

\section{Biểu diễn lời và cảm xúc văn bản tiếng Việt}
\label{section:3.4}
Valence (mức vui--buồn) của một bài hát nằm chủ yếu ở nội dung lời, nên đồ án kết hợp nhiều tín hiệu văn bản độc lập. Điều quan trọng cần phân biệt là ba \textit{loại} mô hình văn bản phục vụ hai mục đích khác nhau, như minh họa trong Hình~\ref{fig:text}: nhóm dự đoán một \textit{giá trị tuyệt đối} (Valence) gồm từ điển và đầu dò trên mô hình ngữ cảnh, còn nhóm \textit{so sánh} hai đoạn lời là mô hình embedding câu.

Tín hiệu từ điển dùng từ điển NRC-VAD \cite{mohammad2018nrcvad} ở bản dịch tiếng Việt, được đối chiếu và bổ sung bằng từ điển cảm xúc bản địa VnEmoLex \cite{doan2022vnemolex} nhằm giảm sai sót dịch máy. Tín hiệu chuyển hóa ngữ cảnh tiếng Việt dùng mô hình ViSoBERT \cite{nguyen2023visobert}, vốn được tiền huấn luyện cho văn bản mạng xã hội tiếng Việt nên hợp với văn phong lời nhạc; mô hình được gắn một đầu dò tuyến tính huấn luyện trên tập cảm xúc UIT-VSMEC \cite{ho2020vsmec}. Tín hiệu liên ngữ dùng mô hình đa ngữ XLM-R \cite{conneau2020xlmr} với đầu dò huấn luyện trên tập EmoBank tiếng Anh \cite{buechel2017emobank}, nhằm tận dụng nguồn nhãn Valence liên tục do con người chấm mà tiếng Việt không có. Các mô hình thay thế như PhoBERT \cite{nguyen2020phobert} đã được cân nhắc; lý do giữ ViSoBERT được trình bày ở Chương~\ref{chapter:SolutionAndContribution}.

Lý do dùng \textit{ba} tín hiệu lời thay vì một là chúng có các kiểu sai khác nhau và bổ trợ cho nhau. Từ điển nhanh và minh bạch nhưng không hiểu ngữ cảnh: nó dễ chấm sai khi gặp phủ định (``không buồn'') hay đảo nghĩa, nên đồ án bổ sung xử lý phủ định và liên từ đối lập cho tín hiệu này. Đầu dò trên mô hình ngữ cảnh (ViSoBERT) hiểu ngữ cảnh và văn phong nhưng phụ thuộc miền tiền huấn luyện --- chính vì văn phong lời nhạc gần với văn bản mạng xã hội mà ViSoBERT được tiền huấn luyện, mô hình này hợp hơn các mô hình huấn luyện trên văn bản trang trọng. Đầu dò liên ngữ (XLM-R trên EmoBank) cho phép \textit{chuyển} nguồn nhãn Valence liên tục do con người chấm từ tiếng Anh sang tiếng Việt --- nguồn nhãn mà tiếng Việt không có. Khi tổ hợp, những lỗi độc lập của ba tín hiệu có xu hướng triệt tiêu lẫn nhau, và bộ tổ hợp EWE (mục~\ref{section:3.5}) tự giảm trọng số cho tín hiệu kém nhất quán với phần còn lại. Bảng~\ref{table:text-signals} tóm tắt bốn tín hiệu văn bản kèm mô hình, vai trò và nguồn nhãn của từng tín hiệu.

\begin{longtable}{|p{2.4cm}|p{3.1cm}|p{2.7cm}|p{3.0cm}|}
\caption{Bốn tín hiệu văn bản: ba tín hiệu dự đoán Valence tuyệt đối và một tín hiệu so sánh tương đồng lời}\label{table:text-signals}\\
\hline
\textbf{Tín hiệu} & \textbf{Mô hình / tài nguyên} & \textbf{Vai trò} & \textbf{Nguồn nhãn} \\
\hline
\endfirsthead
\hline
\textbf{Tín hiệu} & \textbf{Mô hình / tài nguyên} & \textbf{Vai trò} & \textbf{Nguồn nhãn} \\
\hline
\endhead
Từ điển & NRC-VAD-VN + VnEmoLex & Valence tuyệt đối & Từ điển VAD \cite{mohammad2018nrcvad, doan2022vnemolex} \\
\hline
Ngữ cảnh VN & ViSoBERT (đóng băng) + đầu dò & Valence tuyệt đối & UIT-VSMEC \cite{nguyen2023visobert, ho2020vsmec} \\
\hline
Liên ngữ & XLM-R + đầu dò & Valence tuyệt đối & EmoBank (tiếng Anh) \cite{conneau2020xlmr, buechel2017emobank} \\
\hline
Tương đồng lời & multilingual-e5-large & So sánh nội dung lời (cosine) & --- \cite{wang2024e5} \\
\hline
\end{longtable}

Việc chọn tập huấn luyện cho đầu dò ngữ cảnh tiếng Việt cũng là một quyết định có chủ đích về \textit{khớp miền}. Đồ án chọn tập cảm xúc UIT-VSMEC \cite{ho2020vsmec} (cảm xúc trong văn bản mạng xã hội tiếng Việt) thay vì các tập phân tích quan điểm phổ biến hơn như UIT-VSFC \cite{nguyen2018vsfc} (phản hồi sinh viên), vì văn phong và sắc thái cảm xúc của lời nhạc gần với ngôn ngữ mạng xã hội hơn là với văn bản phản hồi trang trọng; một đầu dò huấn luyện trên miền quá xa dễ khiến toàn bộ kho nhạc bị đẩy về một vùng cảm xúc hẹp. Đây là một ví dụ cho nguyên tắc xuyên suốt: với học chuyển giao trên đặc trưng đóng băng, \textit{miền của dữ liệu nhãn phải khớp với miền áp dụng} thì tín hiệu mới đáng tin.

Ngôn ngữ tiếng Việt còn đặt ra những thách thức riêng cho phân tích cảm xúc ở mức từ vựng, củng cố thêm lý do không thể chỉ dựa vào từ điển. Phủ định trong tiếng Việt có thể đứng cách xa từ mang cảm xúc (``chẳng còn thấy vui''); các từ tăng cường (``rất'', ``quá'', ``vô cùng'') và giảm nhẹ (``hơi'', ``khá'') làm thay đổi cường độ cảm xúc; hiện tượng láy (``buồn buồn'') làm dịu sắc thái; và nhiều câu hát mang nghĩa bóng hoặc mỉa mai mà tra từ điển không thể nắm bắt. Đầu dò ngữ cảnh (ViSoBERT) bù đắp phần lớn các trường hợp này nhờ học được biểu diễn phụ thuộc ngữ cảnh, còn đầu dò liên ngữ (XLM-R) mang thêm tri thức cảm xúc từ nguồn nhãn tiếng Anh phong phú. Chính sự bổ trợ giữa ba loại tín hiệu --- từ điển minh bạch nhưng cứng nhắc, mô hình ngữ cảnh linh hoạt nhưng phụ thuộc miền, mô hình liên ngữ giàu nhãn nhưng cần chuyển ngữ --- là lý do tổ hợp nhiều tín hiệu cho kết quả ổn định hơn bất kỳ tín hiệu đơn lẻ nào.

\begin{figure}[tbp]
\centering
\begin{tikzpicture}[font=\small,
   blk/.style={draw, rounded corners, align=center, minimum height=0.8cm, inner sep=3pt},
   src/.style={blk, fill=black!4, text width=2.2cm},
   mid/.style={blk, fill=blue!6, text width=3.2cm},
   outv/.style={blk, fill=green!7, text width=2.0cm},
   arr/.style={-{Stealth[length=2mm]}, thick}]
  \node[src] (lyr) at (0,-2.2) {Lời bài hát};
  % branch 1 lexicon
  \node[mid] (lex) at (4.2,-0.3) {Từ điển NRC-VAD-VN + VnEmoLex};
  % branch 2 contextual probe
  \node[mid] (enc) at (4.2,-2.2) {Encoder ngữ cảnh (ViSoBERT / XLM-R, đóng băng)};
  \node[blk, fill=blue!10, text width=2.0cm] (pb) at (8.4,-2.2) {Đầu dò tuyến tính};
  % branch 3 sentence embedding
  \node[mid] (se) at (4.2,-4.1) {Embedding câu (e5-large)};
  \node[blk, fill=blue!10, text width=2.0cm] (cos) at (8.4,-4.1) {Cosine};
  % outputs
  \node[outv] (val) at (11.7,-1.25) {Valence (tuyệt đối)};
  \node[outv] (simout) at (11.7,-4.1) {Độ tương đồng lời};
  % fan-out từ "Lời bài hát" tới ba nhánh
  \coordinate (fk) at (1.9,-2.2);
  \draw (lyr.east) -- (fk);
  \draw[arr] (fk) |- (lex.west);
  \draw[arr] (fk) -- (enc.west);
  \draw[arr] (fk) |- (se.west);
  % nhánh đi qua đầu dò / cosine
  \draw[arr] (enc.east) -- (pb.west);
  \draw[arr] (se.east) -- (cos.west);
  % gộp về Valence: từ điển vào cạnh trên, đầu dò vào cạnh dưới (cùng đường dọc, không cắt nhau)
  \draw[arr] (lex.east) -| (val.north);
  \draw[arr] (pb.east) -| (val.south);
  \draw[arr] (cos.east) -- (simout.west);
\end{tikzpicture}
\caption{Ba loại tín hiệu văn bản: từ điển và đầu dò ngữ cảnh \textit{dự đoán} Valence tuyệt đối (tổ hợp ở mục~\ref{section:3.5}), còn embedding câu \textit{so sánh} độ tương đồng nội dung lời}
\label{fig:text}
\end{figure}

Các đầu dò cảm xúc nói trên đều là \textit{hồi quy tuyến tính có ràng buộc} (ridge) trên đặc trưng đóng băng, theo Phương trình~\eqref{eq:ridge}: mô hình lớn chỉ đóng vai trò trích đặc trưng $\phi(x)$, còn việc ánh xạ sang giá trị cảm xúc do một lớp tuyến tính đảm nhiệm. Cách này giữ số tham số huấn luyện rất nhỏ, tránh quá khớp khi dữ liệu đích ít và lệch miền.

\begin{equation}\label{eq:ridge}
  \hat{y} = \mathbf{w}^\top \phi(x) + b, \qquad
  \mathbf{w} = \arg\min_{\mathbf{w}} \sum_{i=1}^{n}\big(y_i - \mathbf{w}^\top\phi(x_i)\big)^2 + \lambda \lVert \mathbf{w}\rVert^2
\end{equation}

Đối với chức năng gợi ý bài tương tự, để đo độ tương đồng nội dung lời, đồ án dùng một mô hình embedding câu (sentence embedding) theo kiến trúc Sentence-BERT \cite{reimers2019sbert}. Mô hình embedding câu khác về mục đích so với các đầu dò trên: nó tạo vector để \textit{so sánh} hai đoạn lời bằng cosine, trong khi các đầu dò cảm xúc \textit{dự đoán} một giá trị tuyệt đối. Việc lựa chọn cụ thể giữa một mô hình tiếng Việt chuyên biệt \cite{dangvantuan2024vnembedding} và một mô hình đa ngữ mạnh hơn \cite{wang2024e5} được kiểm chứng bằng thực nghiệm và trình bày ở Chương~\ref{chapter:Evaluation}.

\section{Phương pháp tổ hợp, hậu xử lý và sắp xếp}
\label{section:3.5}
Hệ thống kết hợp nhiều nguồn tín hiệu khác phương thức (âm thanh, lời, màu sắc) nên về bản chất là một bài toán \textit{tổ hợp đa phương thức} (multimodal fusion) \cite{baltrusaitis2019multimodal}; Hình~\ref{fig:fusion} minh họa cách ba phương thức được quy về cùng trục Valence--Arousal. Phần này trình bày các phương pháp tổ hợp, hậu xử lý và sắp xếp tương ứng.

\begin{figure}[tbp]
\centering
\begin{tikzpicture}[font=\small,
   mod/.style={draw, rounded corners, align=center, minimum width=2.7cm, minimum height=1.0cm, fill=black!3},
   fuse/.style={draw, thick, rounded corners, align=center, text width=3.0cm, minimum height=3.8cm, fill=blue!8},
   res/.style={draw, rounded corners, align=center, minimum width=2.4cm, minimum height=1.0cm, fill=green!8},
   arr/.style={-{Stealth[length=2mm]}, thick}]
  \node[mod] (a) at (0,1.6) {Âm thanh\\\footnotesize MuQ};
  \node[mod] (l) at (0,0) {Lời\\\footnotesize e5-large + đầu dò};
  \node[mod] (c) at (0,-1.6) {Màu sắc\\\footnotesize Oklab};
  \node[fuse] (f) at (4.6,0) {Tổ hợp đa phương thức trên trục \textbf{Valence--Arousal}};
  \node[res] (o) at (10.4,0) {Danh sách\\gợi ý};
  \draw[arr] (a.east) -- (f.west|-a); \draw[arr] (l.east) -- (f.west); \draw[arr] (c.east) -- (f.west|-c);
  \draw[arr] (f.east) -- node[above, font=\scriptsize]{truy hồi + xếp hạng} (o.west);
\end{tikzpicture}
\caption{Tổ hợp đa phương thức: âm thanh, lời và màu sắc được quy về cùng trục Valence--Arousal rồi tổ hợp để truy hồi và xếp hạng}
\label{fig:fusion}
\end{figure}

Để tổ hợp nhiều tín hiệu Valence mà không tự đặt trọng số, đồ án dùng bộ ước lượng có trọng số theo độ tin cậy EWE (Evaluator Weighted Estimator) \cite{grimm2005ewe}. Ý tưởng của EWE là: mỗi tín hiệu được cân theo mức tương quan của nó với \textit{đồng thuận chung} của các tín hiệu, nên những tín hiệu nhất quán với phần còn lại được tin hơn (Phương trình~\eqref{eq:ewe}). Cách này không cần nhãn đích nên tránh được tính vòng tròn (không ``học'' để khớp một bộ nhãn có sẵn), đồng thời tự giảm trọng số cho tín hiệu nhiễu.

\begin{equation}\label{eq:ewe}
  \hat{v} = \frac{\sum_{k} w_k\, v_k}{\sum_{k} w_k}, \qquad
  w_k \;\propto\; \operatorname{corr}\!\Big(v_k,\; \tfrac{1}{K-1}\textstyle\sum_{j\neq k} v_j\Big)
\end{equation}

Khi đo tương đồng trên biểu diễn học sâu, các cosine thô bị thiên lệch dị hướng (anisotropy): mọi vector dồn về một nón hẹp khiến cosine giữa hai vector bất kỳ đều cao và gần như mất khả năng phân biệt \cite{ethayarajh2019anisotropy}. Hình~\ref{fig:aniso} minh họa hiện tượng này. Đồ án khắc phục bằng cách trừ trung bình (mean-centering) theo các kết quả đã công bố \cite{mu2018allbutthetop}: trừ đi vector trung bình của toàn kho trước khi tính cosine (Phương trình~\eqref{eq:center}), nhờ đó khôi phục dải động của độ tương đồng.

\begin{equation}\label{eq:center}
  \tilde{\mathbf{x}} = \mathbf{x} - \boldsymbol{\mu}, \qquad
  \boldsymbol{\mu} = \frac{1}{N}\sum_{i=1}^{N}\mathbf{x}_i
\end{equation}

\begin{figure}[tbp]
\centering
\begin{tikzpicture}[font=\small, arr/.style={-{Stealth[length=1.8mm]}, semithick}]
  % left: anisotropic cone (chú thích đặt DƯỚI vòng tròn, không đè mũi tên)
  \begin{scope}[shift={(0,0)}]
    \node[font=\footnotesize\bfseries] at (0,2.5) {Trước (dị hướng)};
    \draw[gray] (0,0) circle (1.8);
    \foreach \ang in {54,60,66,72,78} {\draw[arr, blue!60] (0,0)--(\ang:1.7);}
    \node[font=\scriptsize, align=center] at (0,-2.6) {vector dồn vào nón hẹp\\(cosine $\approx 1$)};
  \end{scope}
  % right: after centering
  \begin{scope}[shift={(6.6,0)}]
    \node[font=\footnotesize\bfseries] at (0,2.5) {Sau khi trừ trung bình};
    \draw[gray] (0,0) circle (1.8);
    \foreach \ang in {20,75,135,190,250,310} {\draw[arr, green!50!black] (0,0)--(\ang:1.7);}
    \node[font=\scriptsize, align=center] at (0,-2.6) {vector trải rộng\\(cosine có nghĩa)};
  \end{scope}
\end{tikzpicture}
\caption{Hiện tượng dị hướng và hiệu ứng trừ trung bình: trước khi trừ, các vector dồn vào một nón hẹp nên cosine gần như bằng nhau; sau khi trừ trung bình, chúng trải rộng và cosine khôi phục khả năng phân biệt}
\label{fig:aniso}
\end{figure}

Để đa dạng hóa kết quả và tránh lặp nghệ sĩ, đồ án dùng độ liên quan biên cực đại MMR \cite{carbonell1998mmr}: ở mỗi bước, hệ thống chọn ứng viên cân bằng giữa \textit{độ liên quan} với truy vấn và \textit{độ khác biệt} so với các bài đã chọn, theo Phương trình~\eqref{eq:mmr}. Cơ chế MMR này được dùng cho gợi ý \textit{bài tương tự}; còn với gợi ý \textit{theo màu}, việc đa dạng hóa nghệ sĩ được lồng trực tiếp vào bước chọn cụm nhất quán (phạt giảm điểm các bài trùng nghệ sĩ) thay vì một bước MMR riêng (chi tiết ở Chương~\ref{chapter:Design}).

\begin{equation}\label{eq:mmr}
  \text{next} = \arg\max_{i \in R \setminus S}\Big[\,\lambda\,\operatorname{rel}(i) - (1-\lambda)\max_{j \in S}\operatorname{sim}(i,j)\,\Big]
\end{equation}

Cuối cùng, để sắp xếp danh sách gợi ý theo màu thành một chuỗi chuyển cảm xúc mượt, đồ án áp dụng nguyên lý đồng cảm (iso-principle) trong trị liệu âm nhạc \cite{altshuler1948iso, starcke2021iso}: chuỗi bài được sắp để tâm trạng thay đổi dần dần dọc theo một đường nối hai điểm cảm xúc, thay vì nhảy đột ngột. Nguyên lý này có gốc rễ trong trị liệu âm nhạc, nơi người trị liệu dẫn dắt tâm trạng người nghe từ trạng thái hiện tại tới trạng thái đích bằng các bước chuyển nhỏ; áp dụng vào danh sách phát hai màu, hệ thống tạo một ``hành trình'' từ tâm trạng của màu A sang màu B.

Cần làm rõ thêm về hai phương pháp ít quen thuộc. EWE \cite{grimm2005ewe} bắt nguồn từ bài toán nhận dạng cảm xúc trong tiếng nói, nơi nhiều người chấm cùng một mẫu và độ tin cậy của mỗi người chấm được ước lượng lặp bằng mức tương quan của họ với trung bình của những người còn lại; ở đây ``người chấm'' được thay bằng ``tín hiệu cảm xúc''. Hiện tượng dị hướng của biểu diễn học sâu \cite{ethayarajh2019anisotropy} đã được quan sát rộng rãi ở các mô hình ngôn ngữ: các vector ngữ cảnh chiếm một hình nón hẹp trong không gian, làm cosine giữa hai vector bất kỳ đều dương và cao. Phương pháp trừ thành phần trội chung (\textit{all-but-the-top}) \cite{mu2018allbutthetop} cho thấy chỉ cần loại bỏ trung bình và một vài thành phần chính là khôi phục được khả năng phân biệt; đồ án dùng phiên bản đơn giản nhất là trừ trung bình, vừa đủ để cải thiện rõ rệt độ nhất quán âm thanh (mục~\ref{subsection:eval-color}).

Ràng buộc ``không tinh chỉnh, chỉ đầu dò tuyến tính'' được củng cố bởi bằng chứng rằng đầu dò tuyến tính trên đặc trưng đóng băng có thể tổng quát hóa tốt hơn tinh chỉnh khi dữ liệu đích lệch miền \cite{kumar2022finetuning}, vốn đúng với trường hợp chuyển từ dữ liệu tiếng Anh sang nhạc Việt.

Về \textit{cấp độ} tổ hợp, trong tổ hợp đa phương thức thường phân biệt tổ hợp sớm (early fusion --- nối các vector đặc trưng thô rồi học chung) và tổ hợp muộn (late fusion --- xử lý từng phương thức riêng rồi hợp nhất ở mức quyết định) \cite{baltrusaitis2019multimodal}. Đồ án chọn một dạng tổ hợp muộn đặc thù: mỗi phương thức (âm thanh, lời, màu) trước hết được quy về \textit{cùng một trục cảm xúc Valence--Arousal có ý nghĩa diễn giải được}, rồi mới hợp nhất. Lựa chọn này có ba lợi ích so với việc nối thẳng các embedding khác chiều và khác thang. Thứ nhất, nó cho phép \textit{diễn giải}: vì điểm hợp nhất nằm trên trục cảm xúc, hệ thống có thể giải thích ``vì sao'' một bài được gợi ý (qua đối tượng \texttt{why} ở Chương~\ref{chapter:Implementation}). Thứ hai, nó tránh để một phương thức có số chiều lớn lấn át khi đo khoảng cách. Thứ ba, nó tách bạch hai loại tín hiệu khác bản chất --- giá trị cảm xúc \textit{tuyệt đối} (V-A) và độ \textit{tương đồng} (cosine embedding) --- để mỗi loại được dùng đúng vai trò trong hàm điểm tổ hợp (Chương~\ref{chapter:Design}).

\section{Công cụ và nền tảng triển khai}
\label{section:3.6}
Phần xử lý dữ liệu, trích đặc trưng và các đầu dò được hiện thực bằng Python với các thư viện học máy phổ biến (NumPy, scikit-learn), các mô hình được tải từ kho mô hình mở. Tầng phục vụ dùng FastAPI vì hỗ trợ tốt kiểu dữ liệu và sinh tài liệu API tự động. Giao diện người dùng được xây dựng dưới dạng ứng dụng trang đơn (Single-Page Application, SPA) bằng React, cho phép tương tác mượt và dựng các hiệu ứng trực quan hóa cảm xúc. Lựa chọn FastAPI cho tầng phục vụ xuất phát từ việc nó hỗ trợ kiểu dữ liệu chặt (giúp xác thực đầu vào ở biên, đáp ứng yêu cầu xử lý lỗi tường minh) và sinh tài liệu API tự động, đồng thời tích hợp tốt với hệ sinh thái Python nơi các đặc trưng đã được tính. Toàn bộ phép so khớp ở pha phục vụ được vector hóa bằng NumPy: thay vì lặp trên từng bài, hệ thống thực hiện một phép nhân ma trận duy nhất giữa vector truy vấn và ma trận đặc trưng của cả kho --- đây là chìa khóa cho thời gian phản hồi mili-giây. Phần liệt kê chi tiết công cụ kèm phiên bản được trình bày ở mục~\ref{section:4.3} của Chương~\ref{chapter:Implementation}.

\section{Ánh xạ yêu cầu sang nền tảng}
\label{section:3.7}
Để thấy rõ mỗi nền tảng phục vụ yêu cầu nào, Bảng~\ref{table:req2found} đối chiếu các yêu cầu (Chương~\ref{chapter:Related_works}) với thành phần lý thuyết/công nghệ đáp ứng. Bảng này cũng cho thấy mọi thành phần đều có nguồn gốc khoa học rõ ràng, đáp ứng yêu cầu NFR01.

\begin{longtable}{|p{5.2cm}|p{8.0cm}|}
\caption{Ánh xạ giữa yêu cầu của hệ thống và nền tảng lý thuyết/công nghệ}\label{table:req2found}\\
\hline
\textbf{Yêu cầu} & \textbf{Nền tảng đáp ứng} \\
\hline
\endfirsthead
\hline
\textbf{Yêu cầu} & \textbf{Nền tảng đáp ứng} \\
\hline
\endhead
Biểu diễn cảm xúc thống nhất (FR01--FR03, FR07) & Mô hình V-A của Russell (mục~\ref{section:3.1}) \\
\hline
Gợi ý theo màu (FR02, FR03) & Oklab + ICEAS + Whiteford (mục~\ref{section:3.2}) \\
\hline
Đo ``nghe giống'' \& suy Arousal (FR01) & SSL âm nhạc MuQ (mục~\ref{section:3.3}) \\
\hline
Suy Valence từ lời (FR07) & Từ điển + đầu dò ViSoBERT/XLM-R (mục~\ref{section:3.4}) \\
\hline
Tổ hợp tín hiệu không tự đặt trọng số (NFR01) & EWE (mục~\ref{section:3.5}) \\
\hline
Kết quả nhất quán âm thanh (FR02) & Trừ trung bình khử dị hướng (mục~\ref{section:3.5}) \\
\hline
Đa dạng hóa \& hành trình mượt (FR03, FR05) & MMR (bài tương tự) / phạt nghệ sĩ trong cụm (theo màu) + iso-principle (mục~\ref{section:3.5}) \\
\hline
\end{longtable}

\section{Phát biểu bài toán hình thức}
\label{section:3.8}
Trên cơ sở các nền tảng đã trình bày, có thể phát biểu hình thức hai bài toán mà hệ thống giải quyết. Gọi $\mathcal{S} = \{s_1, \dots, s_N\}$ là kho nhạc với $N = 5138$ bài. Pha ngoại tuyến gán cho mỗi bài $s_i$ một bộ ba biểu diễn: tọa độ cảm xúc $\mathbf{p}_i \in [0,1]^2$ (Valence--Arousal), vector embedding âm thanh $\mathbf{m}_i \in \mathbb{R}^{1024}$ (MuQ), và vector embedding lời $\mathbf{l}_i \in \mathbb{R}^{1024}$ (e5-large).

\textit{Bài toán 1 (gợi ý bài tương tự).} Cho một bài nguồn $s_q$, tìm tập $k$ bài $\mathcal{R} \subseteq \mathcal{S} \setminus \{s_q\}$ cực đại hóa tổng điểm tương đồng tổ hợp theo Phương trình~\eqref{eq:prob-similar}:
\begin{equation}\label{eq:prob-similar}
  \mathcal{R} = \operatorname*{arg\,top\text{-}k}_{i \,:\, s_i \notin \text{cover}(s_q)} \big[\,0.76\cos(\mathbf{m}_q, \mathbf{m}_i) + 0.16\,\text{sim}(\mathbf{p}_q, \mathbf{p}_i) + 0.08\cos(\mathbf{l}_q, \mathbf{l}_i)\,\big]
\end{equation}
sau đó áp đa dạng hóa MMR trên biểu diễn lời (e5-large) (gián tiếp giảm trùng nghệ sĩ; giới hạn cứng số bài mỗi nghệ sĩ là tùy chọn, mặc định tắt).

\textit{Bài toán 2 (gợi ý theo màu).} Cho một (hoặc hai) màu, ánh xạ thành điểm đích $\mathbf{q}$ trong không gian phân vị V-A của kho nhạc, rồi tìm tập $k$ bài vừa cực đại hóa độ gần cảm xúc (RBF) vừa cực đại hóa độ nhất quán âm thanh theo Phương trình~\eqref{eq:prob-color}:
\begin{equation}\label{eq:prob-color}
  \mathcal{R} = \operatorname*{arg\,top\text{-}k}_{i} \big[\,\alpha\, s_{\text{RBF}}(\mathbf{q}, \mathbf{p}_i) + (1-\alpha)\,\text{coh}(\mathbf{m}_i, \mathcal{R})\,\big], \quad \alpha = 0.55
\end{equation}
trong đó $\text{coh}$ là độ nhất quán âm thanh đo trên biểu diễn MuQ đã trừ trung bình, và kết quả được sắp theo nguyên lý đồng cảm khi có hai màu. Điểm chung của hai bài toán là cùng làm việc trên các biểu diễn $(\mathbf{p}, \mathbf{m}, \mathbf{l})$ đã tính sẵn, nên mọi truy vấn quy về các phép đại số tuyến tính --- nền tảng cho hiệu năng phục vụ ở mức mili-giây (Chương~\ref{chapter:Evaluation}).

Chương này đã trình bày các nền tảng lý thuyết và công nghệ làm cơ sở cho hệ thống, từ mô hình cảm xúc, cơ sở màu sắc, các mô hình biểu diễn âm nhạc và văn bản, đến các phương pháp tổ hợp và sắp xếp, kèm các công thức hình thức hóa cho từng phương pháp. Cách các thành phần này được thiết kế, hiện thực và đánh giá cụ thể cho hai chức năng cốt lõi được trình bày trong chương tiếp theo.

\end{document}
